\section{Energy intervals and the physical seed}\label{sec:seed}
\subsection{Temple's inequality with an absolute threshold}
For a normalized vector $\phi$, let $\mu=\ip\phi{H\phi}$ and
$\sigma^2=\norm{(H-\mu)\phi}^2$. If every nonlowest ordered eigenvalue is at
least $\theta>\mu$, the finite-dimensional Temple inequality is
\begin{equation}\label{eq:temple}
 E_0\ge\mu-\frac{\sigma^2}{\theta-\mu}.
\end{equation}
Indeed, $(E-E_0)(E-\theta)\ge0$ at every eigenvalue, including the lowest one.
Taking its expectation gives
$\sigma^2\ge(\mu-E_0)(\theta-\mu)$.
The input $\theta$ is an \emph{absolute energy threshold}, not a presumed gap.
Proposition~\ref{prop:reference} supplies $\theta=c_A+a$ independently of the
trial. The following positive rational numbers satisfy, throughout $[0,1/10]$,
$d_A\le c_A+a-e_A-m_A$:
\begin{equation}\label{eq:denominators}
 d_{16}=\frac{588758821}{787000000},\qquad
 d_{32}=\frac{568075073}{787000000},\qquad
 d_{64}=\frac{5263369}{7870000}.
\end{equation}
For example, subtracting the endpoint value of
$c_A+a-e_A-m_A$ leaves $(1/10-g)$ times an explicitly nonnegative polynomial;
\lean{ScalarInterval.denominator16_lower} and its two companions prove this
whole-interval statement. There is no endpoint-sampling inference.
Combining \eqref{eq:temple} with Proposition~\ref{prop:trialbounds} gives
\begin{equation}\label{eq:energyinterval}
 e_A-m_A-\varepsilon_A^2/d_A\le E_A\le e_A+m_A
 \quad(A=16,32),\qquad E_{64}\le e_{64}+m_{64}.
\end{equation}
Here $E_A$ denotes the ground energy of its specified box, not an energy
assumed to depend only on area.

\subsection{Separate spatial defects and thermal tails}
Define the two finite ground-energy defects
\[
 d_x^{\rm phys}=E_{8,4}-2E_{4,4},\qquad
 d_y^{\rm phys}=E_{8,8}-2E_{8,4}.
\]
The extensive quadratic and cubic terms cancel exactly in \eqref{eq:shifts}:
$e_{32}-2e_{16}=15g^4/4$ and $e_{64}-2e_{32}=15g^4/2$.
Using \eqref{eq:energyinterval} yields the fixed bounds
\begin{align}\label{eq:dxdy}
 d_x^{\rm phys}&\le\overline d_x
 =\frac{810971762131}{588758821000000}
 =\frac3{8000}+m_{32}+2m_{16}+2\varepsilon_{16}^2/d_{16},\\
 d_y^{\rm phys}&\le\overline d_y
 =\frac{957049999441}{284037536500000}
 =\frac3{4000}+m_{64}+2m_{32}+2\varepsilon_{32}^2/d_{32}.\notag
\end{align}
This correlated cancellation is why separate extensive estimates are not
sufficient: their errors would survive the subtraction. Keeping two different
defects also avoids an unnecessarily severe common-error seed.

At $\beta=16$, apply \eqref{eq:tail} with the actual variational upper
$U_A=e_A+m_A$. Exact positive Taylor sums through degree 50 prove
\[
 e^{16a}>238947,\quad e^{16d_{16}}>157894,\quad
 e^{16d_{32}}>103691,\quad e^{16d_{64}}>44383.
\]
Substitution in the last expression of \eqref{eq:tail} gives
\begin{equation}\label{eq:Qbounds}
 Q_{4,4}<\overline Q_0=\frac{200061}{200000},\quad
 Q_{8,4}<\overline Q_1=\frac{1000927}{1000000},\quad
 Q_{8,8}<\overline Q_2=\frac{1004331}{1000000}.
\end{equation}
The inequalities for the Taylor sums, their application to the real exponential,
and the subsequent rational powers are separate proved steps. Floating-point
exponential evaluation is not used for acceptance.

For a volume $v=(n,m,\beta)$ define the staircase ratios
\begin{equation}\label{eq:staircase}
 A(v)=\frac{Z_{2n,m}(\beta)}{Z_{n,m}(\beta)^2},\quad
 B(v)=\frac{Z_{2n,2m}(\beta)}{Z_{2n,m}(\beta)^2},\quad
 C(v)=\frac{Z_{2n,2m}(2\beta)}{Z_{2n,2m}(\beta)^2}.
\end{equation}
These ratios are not yet assumed to be purities; that is established in
Section~\ref{sec:transfer}. Direct cancellation of the ground-energy factors
and $Q\ge1$ give at $v_0=(4,4,16)$
\[
 A(v_0)\ge\frac{e^{-16\overline d_x}}{\overline Q_0^2},\qquad
 B(v_0)\ge\frac{e^{-16\overline d_y}}{\overline Q_1^2},\qquad
 C(v_0)\ge\frac1{\overline Q_2^2}.
\]
The lower Taylor polynomial $e^{-z}\ge1-z+z^2/2-z^3/6$ for $z\ge0$ now
provides an entirely rational initialization.
\begin{proposition}[Actual initial seed]\label{prop:seed}
For every $g\in[0,1/10]$,
\begin{equation}\label{eq:initial}
 1-A(v_0)\le\frac{11197089}{500000000},\quad
 1-B(v_0)\le\frac{27119173}{500000000},\quad
 1-C(v_0)\le\frac{8606051}{1000000000}.
\end{equation}
\end{proposition}
All five finite inputs needed for \eqref{eq:energyinterval}, and therefore all
three inputs in \eqref{eq:initial}, are proved for the actual Hamiltonian.
The statement is not conditional on a numerical oracle.
At $g=0$ the checks have a useful exact consistency limit:
$Z_{n,m}(\beta)=(1+3e^{-\beta})^{nm}$, $A=B=1$, and
$C=[(1+3e^{-2\beta})/(1+3e^{-\beta})^2]^{4nm}$.
The true gap is then one; the present uniform constant is intentionally much
less sharp.
